% ── Section 2: Background ────────────────────────────────────────────
\section{Background}
\label{sec:bg}

\subsection{Large Language Models and Prompt-Based Interaction}

Modern LLMs such as GPT-4~\cite{achiam2023gpt4} are autoregressive transformer networks trained on broad internet corpora. At inference time they accept a token sequence and produce a continuation; the ``prompt'' is simply the prefix that steers this continuation. In practice, commercial LLM APIs expose a \emph{multi-role message format}: a \texttt{system} message sets the application's persona and policy, followed by alternating \texttt{user} and \texttt{assistant} messages that carry the conversation. All three roles are serialized into a single token stream before the model processes them---meaning the separation between instructions and data exists only by convention, not by any cryptographic or type-level guarantee~\cite{liu2024prompt}. This conflation is the root cause of prompt injection.

The same architecture enables powerful \emph{in-context learning}: by including examples or detailed instructions in the system prompt, developers can customize an LLM's behaviour without retraining. But this flexibility also means that a sufficiently persuasive string in the user or data field can redirect the model just as effectively as the system prompt.

\subsection{Tool-Augmented LLM Applications}

A notable trend since 2023 is augmenting LLMs with the ability to call external tools. The prevailing design pattern interleaves ``thought'' and ``action'' steps: the model reasons about a query and then invokes a function (web search, code execution, database query, email send, etc.) before observing the result~\cite{ruan2024toolemu}. Commercial systems like ChatGPT Plugins and LangChain agents follow this pattern. Retrieval-Augmented Generation (RAG) is a special case: the model's ``tool'' is a vector-search index over documents, and the retrieved text is concatenated into the prompt as context.

Tool augmentation expands the attack surface considerably. Without tools, a compromised LLM can only produce misleading text. With tools, it can read and write files, send messages, execute code, or interact with third-party APIs on the user's behalf. The damage ceiling thus shifts from information manipulation to \emph{autonomous action in the real world}~\cite{ruan2024toolemu}. Ruan et al.\ found that even well-known commercial models triggered unsafe tool calls in 23.9\% of their emulated test scenarios, highlighting the practical risk.

\subsection{Threat Model}
\label{sec:threat}

We define three attacker profiles, each corresponding to one attack family:

\textbf{Direct injection adversary.} The attacker has access to the user-input field and crafts a prompt intended to override the system instructions. This is the simplest threat model: the attacker interacts with the LLM through its normal interface. Examples include ``Ignore all previous instructions'' prefixes and role-playing exploits~\cite{perez2022hackaprompt}.

\textbf{Indirect injection adversary.} The attacker does not interact with the LLM directly. Instead, the attacker plants malicious payloads in data sources that the LLM will retrieve---web pages, email bodies, shared documents, or database records~\cite{greshake2023indirect}. The user, who triggers the retrieval, is typically unaware of the embedded instructions. This model is more realistic for production RAG systems and web-browsing agents.

\textbf{Tool-use adversary.} This attacker targets the LLM's ability to invoke tools. Through either direct or indirect injection, the goal is to cause the model to call a tool with attacker-chosen parameters---for instance, exfiltrating conversation history via an API call or writing malicious files~\cite{zhan2024injecagent}. The attacker may also chain multiple tool calls to achieve a complex objective that no single injection could accomplish.

On the defender side, we assume the system operator controls the system prompt and the set of available tools, but cannot modify the LLM weights (unless explicitly stated, as in training-time defenses). The operator may add pre- or post-processing layers around the model. End users are assumed non-adversarial but their data environment may be compromised.
